一个训练脚本跑了三万步，损失（Loss）从 2.31 掉到 0.087，最近五百步只在小数点后第四位抖动。同事扫了一眼曲线说：“收敛了。”

这句话到底是客观事实，还是主观感觉？

如果抖到第四位算收敛，那第五位呢？再跑三万步会不会突然暴涨？换个随机种子还成不成立？只要“收敛”还停留在“看着差不多稳了”，这就只是一句无法被检验的口号。

两百多年前的数学家也遇到了完全相同的困境，当时那个词叫“趋近”。
算瞬时速度，要求把时间压缩到零；求切线斜率，要求把两点重合；算曲线面积，要求把割片分割无限细——每一处都需要“靠得任意近，却不真的到达”。但多近算“近”？全凭手感。

十九世纪的数学家们给出的解法极其狠辣：不再去讨论“无限接近”这种玄学心理状态，而是把“趋近”改写成一份随时可以被提现的“验收承诺”。

\subsection*{这场“打赌”是怎么进行的？}

想象你和对手在做一场“精度对赌”：
\begin{enumerate}
    \item 对手先提出容差 $\varepsilon$：“我要求输出结果和目标的误差不能超过 $0.001$。”
    \item 你必须给出控制范围 $\delta$：“没问题，只要你把输入的波动控制在半径 $\delta$ 以内，输出就绝对不会超出你的要求。”
\end{enumerate}

只要不管对手提出多么苛刻的 $\varepsilon$，你总能找到对应的 $\delta$ 兑现承诺，我们就说——极限存在。

趋近不再是一种感觉，而是一场随时可以被提现的应答游戏。

\begin{center}
\begin{tikzpicture}[x=1cm,y=1cm,font=\small]
  \node[font=\small\bfseries] at (1.5,4.75) {概念};
  \node[font=\small\bfseries] at (5.0,4.75) {承诺的条款（游戏规则）};
  \node[font=\small\bfseries] at (10.2,4.75) {你的掌控力 ($\delta$)};
  \draw[black!45] (0,4.45) -- (12.4,4.45);
  \draw[black!40,fill=blue!6] (0,3.15) rectangle (12.4,4.35);
  \draw[black!40,fill=blue!10] (0,1.85) rectangle (12.4,3.05);
  \draw[black!40,fill=red!8]  (0,0.55) rectangle (12.4,1.75);
  
  \node at (1.5,3.75) {\textbf{极限}};
  \node at (1.5,2.45) {\textbf{逐点连续}};
  \node at (1.5,1.15) {\textbf{一致连续}};
  
  \node at (5.0,3.75) {看“目标点附近”的趋势，不在乎该点本身};
  \node at (5.0,2.45) {趋势不仅存在，且恰好等于该点真实取值};
  \node at (5.0,1.15) {不仅每个点都连续，控制尺度能通吃整条区间};
  
  \node at (10.2,3.75) {取决于容差 $\varepsilon$ 和目标位置};
  \node at (10.2,2.45) {取决于容差 $\varepsilon$ 和目标位置};
  \node at (10.2,1.15) {\textbf{只取决于 $\varepsilon$，全局通用}};
\end{tikzpicture}

\emph{从“极限”到“一致连续”，本质上只是这场打赌游戏里的“控制力”在不断升级。}
\end{center}

\subsection*{本章四篇的递进分工}

这一章的四篇文章，就是带着你一步步从“凭手感”走向“立条款”：

\begin{itemize}
\item \hyperref[sec:02-Calculus/02-Limits-and-Continuity/01]{\textbf{“为什么需要极限”}}：区分“在某点的值”与“在某点附近的趋势”。眼睛看到的图像和采样点都会骗人，只有趋势不会。
\item \hyperref[sec:02-Calculus/02-Limits-and-Continuity/02]{\textbf{“把趋近写成严格承诺”}}：建立 $\varepsilon$-$\delta$ 的应答机制。这一篇的核心不是去凑具体的数学计算，而是读懂谁先出手、谁依赖谁。
\item \hyperref[sec:02-Calculus/02-Limits-and-Continuity/03]{\textbf{“极限计算与夹逼”}}：教你如何利用运算法则与夹逼定理拆解承诺，不必每次都回到底层硬推 $\delta$。
\item \hyperref[sec:02-Calculus/02-Limits-and-Continuity/04]{\textbf{“连续性与整体保证”}}：把趋势锁定到函数自身取值上。你将看到，“局部有保证”如何一步步升级为“整段区间上的整体性质”（如介值定理与最值定理）。
\end{itemize}

\textit{提示：第一次读到 $\varepsilon$-$\delta$ 的严格定义时，切记不要死扣具体的算术凑数。先牢记这句话——“你报出你的误差容忍度，我给出我的控制窗口”。抓住了这个游戏规则，后面所有的抽象符号都会顺理成章。}